\section{Experimental Setup}

\subsection{Scenario}

The main evaluation uses a hard re-entry sequence with two people. The selected target is the black-shirt person. The sequence includes crossing, occlusion, identity ambiguity, and re-entry. These conditions stress the ability of the system to preserve the selected target rather than simply maintain any person track.

\subsection{Compared Trackers}

Three base tracker settings are evaluated:

\begin{itemize}
    \item ByteTrack;
    \item OCSORT;
    \item DeepSORT-MARS.
\end{itemize}

For each tracker, the raw selected target is compared against the TIM-MARS selected-target output when available.

\subsection{Metrics}

The evaluation separates selected-target output into three main categories:

\begin{itemize}
    \item correct: the output corresponds to the true selected person;
    \item wrong: the output corresponds to another person;
    \item lost: no valid selected-target output is available while the selected person is visible.
\end{itemize}

The primary ratios are computed over the visible target duration:

\begin{equation}
r_\mathrm{correct} = \frac{t_\mathrm{correct}}{t_\mathrm{visible}},
\end{equation}

\begin{equation}
r_\mathrm{wrong} = \frac{t_\mathrm{wrong}}{t_\mathrm{visible}},
\end{equation}

\begin{equation}
r_\mathrm{lost} = \frac{t_\mathrm{lost}}{t_\mathrm{visible}}.
\end{equation}

The most important safety metric is wrong-target duration.

\subsection{Runtime Measurement Scope}

The runtime comparison isolates the post-detection tracking and selected-target memory stack. The source bag replays \texttt{/camera/image\_raw} and \texttt{/detections}, while Hailo inference is excluded. This makes the comparison suitable for measuring tracker and TIM-MARS overhead under identical detector outputs, but it should not be interpreted as full onboard perception-pipeline runtime.
